\section{Distribute}
To distribute the rules from the central database to the sensors in a format
that \snort{} can understand, three major things happen. Each of these
is covered in the following three subsections.

\subsection{Initiating the Distribute}

Before a synchronization can start, the sensors  need to know that they
should synchronize. There are mainly two things that can cause a synchronization
to occur:

\begin{enumerate}
  \item Automatic synchronization based on time
  \item Manually invoked trough the web interface on the \snowman{}-server
\end{enumerate}

In the first case, the client itself knows that it should start to synchronize
and therefore starts the authentication/synchronization procedure with the
server. 

In the latter case, the client does not initially know that anything happened.
The central server has to contact the client via the ``snowmanclientd''
daemon, telling it to initiate a synchronization.

\subsection{Synchronization}

When a synchronization is initiated the client connects to the
``snowmand'' daemon on the sentral server and authenticates with it. If the
authentication is successful, the server knows that the sensor is whom it
claims to be and it also knows which rules the sensor should have installed.
we can then begin the synchronization.

The first thing that is done is to retrieve all the RuleClass objects and
update the local cache on the client to correspond to the retrieved list from
the server.
Changed objects will be updated, new objects inserted into the local database and obsolete objects
will be deleted. Next lists of Generators, RuleReferenceTypes and RuleSets is
retrieved and the local database is updated in a simmilar manner as with the
RuleClass objects.

For the synchronization of the rules themselves, which is the main bulk of the
data needed to be synchronized, a more efficient approach is used. First a list
of the signature-id's and revisions of the rules that should be present on the
sensor is retrieved from the server. If rules in the local cache is not
mentioned in this list or this list indicates that a new version of this rule is available, the rule is
deleted from the local cache. Any tuning parameters belonging to this rule
is also removed.

Next, the client starts to request the rules that it currently is missing from
the server. To speed things up, this happens in bulks of 250 rules per request.
This way, no single requests gets too big and time-consuming. For the rules
requested, tuning parameters like filters and suppressions is supplied if the
rules have any such applied on the given sensor. References is also passed along
with the rule.

All the recieved rules are then saved to the local database and the
synchronization is complete. The final step of generating \snort{}
configuration files is then started.

\subsection{Generating Snort Configuration}
The generation of the \snort{} configuration files is a rather simple process.
First, the old configuration is cleared out and new files are generated based
on the content of the local database.

The naming convention of the files is similar to the one SourceFire uses with
their rulesets and the format of the files are exactly the same. This is
to simplify the visual inspection of the rulefiles if that is desirable. Snort
in and of itself does not really care if the configuration files are in a
certain order, as long as the content can be recognized as \snort{} configuration.

In addition to the various files created by \snowman{}, one certain file is
guaranteed to always be generated. That is the ``snowman-includes.cfg'' file,
which includes all the files generated. If \snort{} includes that file again,
all the files generated by \snowman{} will be included for \snort{}.

Finally, when the configuration has been generated, the snowman-client sends a
SIGHUP signal to the \snort{} process, which makes \snort{} reload its
configuration and start using its new rulesets.
